% TOP_PROBLEM: {"id": 3066, "title": "Equality in a matroidal circumference bound", "category_id": 2, "category": {"id": 2, "name": "combinatorics", "display_name": "Combinatorics", "description": "Counting problems, graph theory, discrete structures.", "slug": "combinatorics", "order_index": 2, "created_at": "2026-07-31T15:26:25.671Z"}, "status": "open"}
% FIRST_POSTED: null
% ENTRY_KIND: statement_only
% Imported from: batch01.tex; UnsolvedMath ID 3066
\documentclass[11pt]{book}
\usepackage{fontspec}
\setmainfont{Latin Modern Roman}
\setsansfont{Latin Modern Sans}
\usepackage{amsmath,amssymb,mathtools}
\usepackage{unicode-math}
\setmathfont{Latin Modern Math}
\usepackage{xurl}
\usepackage[hidelinks]{hyperref}
\newcommand{\sref}[2]{\hyperlink{source:#1:#2}{[S#2]}}
\newcommand{\cataloguescope}[1]{\paragraph{Mathematical question.}#1}
\begin{document}
% UnsolvedMath ID 3066; source catalogue code OPG-692
\setcounter{section}{3065}
\section{Equality in a matroidal circumference bound}\label{3066}
{\small\sffamily UnsolvedMath ID 3066 · Combinatorics}
\subsection{Definitions and mathematical statement}

\paragraph{Shared notation.}$\mathbb N=\{1,2,\ldots\}$, and $\mathbb Z,\mathbb Q,\mathbb R,\mathbb C$ denote the integers, rationals, reals and complex numbers. Any other conventions are specified locally.


A finite matroid $M$ on ground set $E$ is specified by a family of independent subsets containing $\varnothing$, closed under taking subsets and satisfying the augmentation axiom. Its rank $r_M(A)$ is the largest independent-set size within $A$. Circuits are minimal dependent sets, and $c(M)$ is the maximum circuit size. The dual matroid $M^*$ has bases complementary to bases of $M$. A $k$-separation for $k\in\{1,2\}$ is a partition $E=A\mathbin{\dot\cup}B$ with $|A|,|B|\geq k$ and $r_M(A)+r_M(B)-r_M(E)<k$. Here $3$-connected means $|E|\geq4$ and no one- or two-separation. The affine binary cube $\operatorname{AG}(3,2)$ is the rank-four binary matroid represented by the eight columns $(1,x_1,x_2,x_3)^T$ with $(x_1,x_2,x_3)\in\mathbb F_2^3$. A graphic matroid is the cycle matroid of a finite graph.
\paragraph{Statement recorded in the supplied source.}
Is every $3$-connected matroid satisfying
\[|E|=\frac{c(M)c(M^*)}{2}\]
isomorphic to $\operatorname{AG}(3,2)$? The associated variant asks whether all connected non-graphic equality cases, apart from $\operatorname{AG}(3,2)$ itself, are obtained from it by replacing every element with a parallel class of the same positive size. \sref{3066}{2}
\subsection{Short English statement}
Is the binary affine cube the unique three-connected matroid attaining equality in the circuit--cocircuit size bound, and are its uniform parallel extensions the only non-graphic connected equality variants?
\subsection{Sources}
\begin{description}
\item[\hypertarget{source:3066:1}{S1}] ULAM.AI and the UnsolvedMath Contributors, \emph{UnsolvedMath: A Curated Collection of Open Mathematics Problems}, record 3066 (OPG-692). CC BY 4.0. \url{https://huggingface.co/datasets/ulamai/UnsolvedMath}.
\item[\hypertarget{source:3066:2}{S2}] Oxley, James; Royle, Gordon, Equality in a matroidal circumference bound, Open Problem Garden, node 692. Complete problem, discussion and bibliography (including mathematical image alt text). \url{https://www.openproblemgarden.org/op/equality_in_a_matroidal_circumference_bound}.
\end{description}
\label{end:3066}
\end{document}
